% =============================================================================
% Project : THE_PHILOSOPHY_OF_VIABILITY
% Book    : Book I — Ontology of Viability
% Book ID : VS.B01
%
% File    : 2_MAINMATTER/PART_I_IDENTITY/chapter_01_identity/sections/section_10_scope_and_limits_of_identity.tex
%
% Section / Chapter:
%   The Scope and Limits of Identity
%
% Purpose :
%   To delimit the ontological claims established in the chapter,
%   distinguish foundational commitments from argued conclusions
%   and provisional constructions, and identify the principal
%   unresolved problems concerning identity.
%
% =============================================================================

\section{The Scope and Limits of Identity}
\label{sec:scope-and-limits-of-identity}

\subsection{The Need to Define the Scope of Ontological Claims}

An ontological investigation must establish not only the
distinctions through which a problem can be examined, but also
the limits of the claims that those distinctions support.

Without an explicit account of scope, conclusions may be
extended beyond the conditions under which they have been
established. A conceptual distinction may be mistaken for an
ontological separation, a possible relation may be treated
as a necessary one, or a partial explanation may be presented
as a universal principle.

The preceding sections have developed a framework for
investigating identity through its relations to being,
determination, unity, distinction, individuation, numerical
identity, and signature. They have also established a
constraint against reducing identity to any one of these
dimensions.

These conclusions clarify the conceptual and ontological
questions that an investigation of identity must address.
They do not, however, constitute a complete account of every
condition under which identity is constituted, maintained,
transformed, or manifested.

The purpose of this section is therefore to delimit the
scope of the present account, clarify the status of its
principal claims, and distinguish what has been established
from what remains provisional or unresolved.

The limits identified here concern the extent of the
arguments developed in this chapter. They must not be
confused with limits imposed upon existence or with claims
about what can or cannot be established by subsequent
ontological inquiry.

\subsection{The Ontological Scope of Identity}

The subject of this chapter is the ontological character
of identity.

Its central question concerns what makes an existent the
existent that it is. This question concerns the existent
itself and does not depend upon an act through which it is
identified, classified, named, or described.

Identity is therefore not treated as a product of
identification or as a determination introduced merely
through an external act. Its ontological status belongs
to the existent rather than to an activity directed toward it.

Within this scope, the chapter has examined the following
relations and distinctions:

\begin{enumerate}
    \item The ontological relation between being and identity.
    \item The relation between identity and determination.
    \item The distinction between identity and unity.
    \item The relations among identity, distinction, and difference.
    \item The relation between identity and individuation.
    \item The distinction between identity and numerical identity.
    \item The relation between identity and signature.
    \item The non-reduction of identity to other related
    ontological dimensions.
\end{enumerate}

These examinations establish a conceptual framework for
the further investigation of identity. They do not establish
that every existent has the same constitutive conditions,
that every dimension examined is relevant in every case,
or that the relations identified operate uniformly across
all possible kinds of existents.

The scope of the chapter is thus general in its ontological
question, but limited in the extent of the constitutive
principles it has established.

\subsection{Identity and Identification}

An important boundary of the present investigation concerns
the distinction between identity and identification.

An existent does not acquire its identity because an act
of identification occurs. Nor does the absence of such an
act imply the absence of identity.

The ontological character of an existent is not constituted
by its being named, classified, distinguished, or recognized
by another existent. Such activities may concern an
existent, but they do not provide the ontological ground
of its identity.

This commitment does not require a theory of identification,
nor does it determine the conditions under which identity
may become accessible to another existent. Those questions
concern additional relations and processes that are not
the subject of the present chapter.

The distinction is more limited and more fundamental:
the identity of an existent must not be made ontologically
dependent upon the occurrence of an identifying activity.

Accordingly, the present account treats identity as an
ontological question independently of whether identification
takes place. This does not entail that every possible
relation between an existent and another existent is
irrelevant to identity. It excludes only the assumption
that an act of identification is itself the ground from
which identity derives.

\subsection{The Limits of a Universal Definition}

The analysis developed in this chapter does not establish
a single universal principle that exhaustively determines
identity in every possible existent.

A principle may explain a particular dimension of identity
without explaining its complete ontological character.
An account of individuation, for example, may address the
constitution of an existent as an individual. An account
of invariance may address the preservation of selected
features or relations under specified transformations.
Neither explanatory function, by itself, establishes
a universal definition of identity.

The same limitation applies to principles based exclusively
on material continuity, structural organization, relational
stability, numerical identity, or the preservation of
a signature.

The present investigation therefore does not claim that
any one of these principles provides a universally
necessary and sufficient criterion of identity.

This limitation must be interpreted carefully. The absence
of an established universal criterion does not imply that
identity has no ontological conditions, nor does it imply
that such conditions are arbitrary. It means that their
general form, necessity, sufficiency, and domain of
applicability have not been established by the arguments
developed in this chapter.

Whether a more comprehensive principle can be formulated
remains an open problem.

The possibility of a universal account must therefore
be distinguished from the claim that such an account has
already been established. The former remains available
for investigation; the latter would exceed the present
argument.

\subsection{The Limits of Claims Concerning Persistence and Transformation}

Identity raises the question of how an existent remains
the existent that it is under different conditions and
through transformation.

The distinctions developed in this chapter prevent change
from being treated as automatically equivalent to the
loss of identity. They do not, however, establish a
universal criterion by which every transformation can
be classified as compatible or incompatible with the
persistence of an existent.

Material, structural, relational, and other determinations
may change in different ways and with different consequences.
Their significance cannot be assumed to be identical
across all kinds of existents.

Accordingly, the following distinctions must be maintained:

\begin{enumerate}
    \item Change does not necessarily entail the loss
    of identity.

    \item The continuity of selected determinations
    does not necessarily establish the persistence
    of identity.

    \item The preservation of a structure, signature,
    or invariant does not, by itself, establish
    numerical identity.

    \item The alteration or replacement of constituents
    does not, by itself, establish the destruction
    of an existent.

    \item Numerical identity across different conditions
    does not provide a complete explanation of the
    ontological character of the existent.

    \item The persistence of an existent and the
    preservation of every determination associated
    with it are not equivalent claims.
\end{enumerate}

These distinctions do not determine which transformations
preserve identity or which conditions are necessary for
persistence. They prevent such determinations from being
inferred solely from the presence or absence of selected
features.

A more comprehensive account must examine the relations
among identity, transformation, continuity, discontinuity,
replacement, and numerical identity in relation to the
kinds of existents and conditions under consideration.

The present chapter establishes the need for this
investigation, not its final criteria.

\subsection{The Limits of Claims Concerning Multiple Identities}

The present chapter does not establish whether one
existent can involve more than one identity.

The unity of an existent does not necessarily demonstrate
that its identity is singular in every relevant
ontological sense. Conversely, internal plurality,
compositional complexity, or structural differentiation
does not establish the existence of multiple identities.

The existence of multiple signatures, configurations,
determinations, or manifestations likewise does not entail
the existence of multiple identities. Such multiplicity
may concern different aspects or conditions of one
existent rather than distinct identity determinations.

The question requires a more precise account of the
relations among identity, unity, individuation, and the
ontological constitution of an existent. In particular,
it remains necessary to determine whether distinct
identity determinations could belong to one and the same
existent without implying the existence of multiple
individuals.

This inquiry must distinguish the possible multiplicity
of identity determinations from the numerical multiplicity
of individuals. The former, if established, would not
automatically entail the latter.

Neither the universal singularity nor the possible
plurality of identity can therefore be asserted on the
basis of the distinctions established so far.

The question remains open for subsequent ontological
investigation.

\subsection{The Limits of Universality Across Domains of Existence}

The conceptual framework developed in this chapter is
intended to support a general ontological investigation.
Nevertheless, its applicability across different domains
of existence does not establish that all identity-related
conditions operate identically in every domain.

A general ontological concept may retain its conceptual
function while the determinations through which it is
constituted or manifested vary according to the nature
of the existent.

Material constitution, structural organization, and
relational determination, for example, may have different
forms or degrees of relevance in different domains.
Such variation does not necessarily imply that identity
itself has an entirely different meaning in each domain.

Conversely, conceptual uniformity does not establish
that the ontological conditions of identity are uniform
throughout existence.

The present chapter therefore establishes neither the
complete uniformity nor the fundamental heterogeneity
of identity across all domains. It provides a general
conceptual framework whose applicability to particular
domains requires further examination.

The investigation of different realms of being must
determine how identity relates to their particular
ontological conditions without presupposing either
universal uniformity or absolute discontinuity.

The term ``domain'' is used here to designate a scope
within which particular ontological conditions are
examined. Its use does not establish a definitive
classification or hierarchy of domains, nor does it
presuppose that the proposed organization of the ontology
of viability is exhaustive.

\subsection{The Limits of the Relationship Between Identity and Manifestation}

The relation between identity and manifestation remains
only partially determined within the present investigation.

The chapter has established that identity cannot be
reduced to manifestation and that a particular
manifestation cannot automatically be treated as an
exhaustive account of the identity of an existent.

This does not establish that manifestation is merely
external to identity, or that the identity of an existent
remains unaffected by every possible condition of
manifestation.

A manifestation may involve particular determinations
of an existent under conditions associated with a
specific ontological context or domain. Whether such
conditions participate in the constitution, realization,
or transformation of identity requires further analysis.

The present chapter does not determine the complete
ontological status of manifestation or expression, nor
does it establish the relations among the domains in
which manifestation occurs.

These matters must be examined in the broader
investigation of manifestation and the realms of being.
Their relationship to identity cannot be settled merely
by the conceptual distinctions established here.

\subsection{The Limits of the Present Conceptual Framework}

The distinctions developed in this chapter provide
necessary conceptual conditions for a more comprehensive
ontology of identity. They do not yet provide a complete
account of its constitutive principles.

Several questions remain unresolved:

\begin{enumerate}
    \item What are the necessary and sufficient ontological
    conditions of identity?

    \item Can a general principle of identity be established
    across different domains of existence?

    \item How can constitutive determinations be distinguished
    from non-constitutive determinations in a general
    ontological account?

    \item Under what conditions can identity persist
    through substantial transformation?

    \item What is the complete relationship between
    identity, individuation, and numerical identity?

    \item Can one existent involve multiple identities
    without implying multiple individuals?

    \item What is the precise ontological relationship
    between identity and manifestation?

    \item To what extent are structure, composition,
    material constitution, and relational determination
    constitutive of identity?
\end{enumerate}

These questions identify the limits of the present account.
Their unresolved status does not, by itself, constitute
a contradiction within the framework or provide evidence
against its foundational commitments.

An unresolved question must not be treated as an
established conclusion. Equally, the absence of a
demonstrated principle must not be interpreted as proof
that such a principle is impossible.

The distinction between what has been established and
what remains open is therefore essential to the continued
development of the ontology.

\subsection{The Status of the Claims Developed in This Chapter}

The claims developed in this chapter must be distinguished
according to their argumentative and justificatory status.
Not all claims have the same function, and not all possess
the same degree or kind of support.

The following commitments provide the foundational
orientation of the present investigation:

\begin{enumerate}
    \item Being and identity are ontologically inseparable
    and analytically distinguishable.

    \item Neither being nor identity has ontological
    priority over the other.

    \item Identity is not an additional property attached
    to an identity-free existent.

    \item The ontological status of identity does not
    depend upon an act of identification.
\end{enumerate}

These commitments define the starting orientation of
the chapter. Their foundational role must not be confused
with the claim that every implication of them has already
been demonstrated.

The following distinctions have been developed through
the conceptual and argumentative analysis of the chapter:

\begin{enumerate}
    \item Identity is not conceptually equivalent to
    determination, unity, distinction, individuation,
    numerical identity, or signature.

    \item Conceptual non-equivalence does not establish
    ontological separation.

    \item Non-reduction does not imply ontological isolation.

    \item No single related ontological dimension has
    been established as an exhaustive definition
    of identity.
\end{enumerate}

These are conclusions of the chapter's analysis. Their
scope is limited to the arguments supporting them and
does not extend automatically to every possible
ontological relation among the dimensions concerned.

The following matters remain provisional or require
further ontological justification:

\begin{enumerate}
    \item The precise constitutive roles of material,
    structural, and relational determinations.

    \item The conditions under which particular
    transformations preserve or alter identity.

    \item The applicability of general identity principles
    across different domains of existence.

    \item The precise relation between identity and
    manifestation.
\end{enumerate}

The following remain open ontological problems:

\begin{enumerate}
    \item Whether a universal principle of identity
    can be established.

    \item Whether one existent can involve multiple
    identities.

    \item What constitutes a complete account of the
    necessary and sufficient conditions of identity.

    \item How identity, individuation, continuity,
    and numerical identity are related in a
    comprehensive ontology.
\end{enumerate}

This classification is local to the present chapter.
It is not a replacement for the broader system of
epistemic and justificatory statuses adopted by the
project. In particular, foundational commitments,
stipulated or defined claims, argued or derived claims,
conditionally supported claims, provisional constructions,
hypotheses, and unresolved problems remain distinguishable
within that broader system.

The classification above identifies the principal
groups of claims relevant to the present chapter; it
does not assign every individual proposition a final
status within the complete ontology.

Nor is this classification intended to establish a
permanent boundary around the theory. It clarifies
the current status of its claims and allows subsequent
arguments to develop, qualify, or revise them without
confusing foundational commitments with unresolved
possibilities.

\subsection{Ontological Restraint and the Development of the Theory}

The limits established in this section concern the
extent of the ontological claims justified by the
present investigation. They are not limitations
imposed upon existence itself.

Ontological restraint requires that the scope of every
claim remain proportionate to the argument supporting it.

A conceptual distinction may be established without
resolving every ontological relation between the
concepts involved. A possible condition may be
identified without being affirmed as universally
necessary or actual. A provisional construction may
organize further inquiry without being treated as a
completed ontological foundation.

This restraint is not a suspension of ontological
inquiry. It is a condition for its systematic development.

It also requires that the absence of a demonstrated
explanation be distinguished from the absence of an
ontological relation. The fact that a constitutive
relation has not yet been established does not show
that no such relation exists. Likewise, identifying
a possible relation does not establish that it is
necessary, universal, or sufficient.

The ontology of viability must therefore preserve
the distinctions among foundational commitments,
stipulated or defined claims, argued or derived claims,
conditionally supported claims, provisional constructions,
hypotheses, and open problems throughout its subsequent
development.

Such distinctions allow the theory to develop without
prematurely universalizing partial results or treating
unresolved questions as though they had already been
answered.

\subsection{Concluding Formulation}

The scope of the present investigation is the
ontological character of identity as inseparable from
the being of an existent. Its limits concern the extent
to which the constitutive conditions, principles, and
relations of identity have been established by the
arguments developed in this chapter.

The argument of this chapter supports the following
formulation:

\begin{quote}
Identity is an ontological dimension inseparable from
the being of an existent, but its complete constitution
cannot be inferred from any single related determination,
condition, or relation. The distinctions developed in
this chapter provide the conceptual conditions for
investigating identity without prematurely asserting
a universal criterion, a complete theory of persistence,
the singularity or plurality of identities, or an
exhaustive account of the relationship between identity
and manifestation.
\end{quote}

The present account establishes a foundational framework
for the ontology of identity, not a complete theory of
every condition under which identity is constituted,
maintained, transformed, or manifested.

Its unresolved questions define the direction of further
ontological inquiry rather than a deficiency in the
reality being investigated.

The chapter can therefore conclude with a distinction
between the ontological problem itself and the limits
of its present formulation: the identity of an existent
is not exhausted by the concepts through which it has
been examined, while the relations among those concepts
remain necessary subjects of further investigation.

The subsequent synthesis must preserve both sides of
this commitment: the ontological inseparability of being
and identity, and the need to investigate identity
without reducing it to any single related determination.

% =============================================================================
% END OF SECTION 10 OF CHAPTER 01
% =============================================================================